\section{符号说明}

\begin{table}[H]
  \centering
  \caption{主要符号说明}
  \begin{tabular}{clc}
    \toprule
    \textbf{符号} & \textbf{含义} & \textbf{单位} \\
    \midrule
    $S_1,S_2$                 & 第一、第二检测点的位置矢量 & m \\
    $\theta_1,\theta_2$       & 两个检测点上的示向度（读数） & $^\circ$ \\
    $\varepsilon$             & 测向误差上界，$\varepsilon=1^\circ$ & $^\circ$ \\
    $G,\;d=|S_1G|$            & 干扰源的位置与它到 $S_1$ 的距离 & m \\
    $r_1,r_2$                 & 源到 $S_1$、$S_2$ 的距离 & m \\
    $\gamma$                  & 源处两条示向度的交会角 & $^\circ$ \\
    $\delta_1,\delta_2$       & 两次测向的误差 & $^\circ$ \\
    $\rho_{\mathrm{rec}}$     & 源的接收半径，$\rho_{\mathrm{rec}}\in[1000,1500]$ & m \\
    $C$                       & 源的不确定集（窄楔形） & --- \\
    $D(R)$                    & 以原点为心、半径 $R=1800$ m 的圆域 & --- \\
    $W(p,\psi)$               & 顶点 $p$、中心方向 $\psi$、半张角 $\varepsilon$ 的楔形 & --- \\
    $R(S_1,S_2)$              & 两条示向度围成的定位区域（楔形交） & --- \\
    $\mathrm{diam}(\cdot)$    & 平面集合的直径（最大点距） & m \\
    $J(S_2)$                  & 最坏情况定位直径（代价函数） & m \\
    $J^*,\;S_2^*$             & 最优值与其对应的第二检测点 & m \\
    $F$                       & 可行域（可测性透镜） & --- \\
    $\mathcal{F}(S_2)$        & 适合度 $=J^*/J(S_2)\in(0,1]$ & --- \\
    $\eta$                    & 候选区域阈值参数（取 $10\%$） & --- \\
    $u(\psi)$                 & 单位方向矢量 $(\cos\psi,\sin\psi)$ & --- \\
    $R_1,R_2$                 & 文献判据中的两站测距（记为 $d,r_2$） & m \\
    $\sigma,\varepsilon_{\mathrm{ang}}$ & 文献口径中的测向误差标准差 / 角度 & $^\circ$ \\
    $H$                       & 两站纯方位测向的 Fisher 信息矩阵 & --- \\
    $a$                       & CRLB 误差椭圆主半轴 & m \\
    $S_{\mathrm{ellipse}}$    & CRLB 误差椭圆面积 & m$^2$ \\
    $\rho_{\mathrm{MEC}}$     & 几何稀释式给出的最小外接圆半径 & m \\
    $T_{\max}$                & 单次测向耗时 $5$ s（本问不进入代价） & s \\
    \bottomrule
  \end{tabular}
\end{table}

其余符号在正文首次出现处说明。所有常量、网格步长与结果数值均由确定性程序给出，
$\eta$ 为可调参数（命令行 \texttt{--eta}）。
